% Part: first-order-logic
% Chapter: tableaux
% Section: provability-consistency

\documentclass[../../../include/open-logic-section]{subfiles}

\begin{document}

\iftag{FOL}
      {\olfileid[hi]{fol}{tab}{prv}}
      {\olfileid[hi]{pl}{tab}{prv}}

\olsection{\usetoken{S}{derivability} और संगतता}

अब हम !!{derivability} संबंध के अनेक गुण स्थापित करेंगे। ये अपने-आप में
रोचक हैं, पर इनमें से प्रत्येक पूर्णता प्रमेय के प्रमाण में भूमिका निभाएगा।

\begin{prop}\ollabel{prop:provability-contr}
  यदि $\Gamma \Proves !A$ और $\Gamma \cup \{!A\}$ असंगत है, तो
  $\Gamma$ असंगत है।
\end{prop}

\begin{proof}
ऐसे परिमित $\Gamma_0 = \{!B_1, \dots, !B_n\}$ और $\Gamma_1
  =\{!C_1, \dots, !C_n\} \subseteq \Gamma$ हैं कि
  \begin{align*}
    \{\sFmla{\False}{!A}, &
    \sFmla{\True}{!B_1}, \dots, \sFmla{\True}{!B_n}\} \\
    \{\sFmla{\True}{!A}, &
    \sFmla{\True}{!C_1}, \dots, \sFmla{\True}{!C_m}\}
  \end{align*}
  के लिए बंद !!{tableau}s हैं। $!A$ पर \Cut{} नियम लगाकर इन्हें एक बंद
  !!{tableau} में जोड़ा जा सकता है, जो दिखाता है कि $\Gamma_0
  \cup \Gamma_1$ असंगत है। चूँकि $\Gamma_0
  \subseteq \Gamma$ और $\Gamma_1 \subseteq \Gamma$, इसलिए $\Gamma_0 \cup
  \Gamma_1 \subseteq \Gamma$;
  अतः $\Gamma$ असंगत है।
\end{proof}

\begin{prop}
\ollabel{prop:prov-incons}
$\Gamma \Proves !A$ तभी और केवल तभी जब $\Gamma \cup \{\lnot !A\}$ असंगत हो।
\end{prop}

\begin{proof}
पहले मान लें $\Gamma \Proves !A$; अर्थात् निम्न के लिए कोई बंद
!!{tableau} है:
\[
\{\sFmla{\False}{!A},
\sFmla{\True}{!B_1}, \dots, \sFmla{\True}{!B_n}\}
\]
$\TRule{\True}{\lnot}$ नियम का प्रयोग करके इसे निम्न के लिए बंद
!!{tableau} में बदला जा सकता है:
\[
\{\sFmla{\True}{\lnot !A},
\sFmla{\True}{!B_1}, \dots, \sFmla{\True}{!B_n}\}.
\]

दूसरी ओर, यदि बाद वाले के लिए कोई बंद !!{tableau} है, तो पहली मान्यता
$\sFmla{\True}{\lnot !A}$ और उस पर \TRule{\True}{\lnot} लगाने से प्राप्त
हर सूत्र को हटाकर तथा मान्यता $\sFmla{\False}{!A}$ जोड़कर हम उसे पहले वाले
का बंद !!{tableau} बना सकते हैं। यदि कोई शाखा पहले इसलिए बंद थी कि उसमें
$\sFmla{\True}{\lnot !A}$ पर \TRule{\True}{\lnot} लगाने का निष्कर्ष,
अर्थात् $\sFmla{\False}{!A}$, था, तो नए !!{tableau} में संगत शाखा भी बंद
है। यदि पुराने सत्य-वृक्ष में कोई शाखा इसलिए बंद थी कि उसमें मान्यता
$\sFmla{\True}{\lnot !A}$ के साथ $\sFmla{\False}{\lnot
  !A}$ भी था, तो $\TRule{\False}{\lnot}$ को
$\sFmla{\False}{\lnot !A}$ पर लगाकर
$\sFmla{\True}{!A}$ प्राप्त करें। चूँकि हमने
$\sFmla{\False}{!A}$ को मान्यता के रूप में जोड़ा है, इससे शाखा बंद हो जाती
है।
\end{proof}

\begin{prob}
सिद्ध कीजिए कि $\Gamma \Proves \lnot !A$ तभी और केवल तभी जब
$\Gamma \cup \{!A\}$ असंगत हो।
\end{prob}

\begin{prop}\ollabel{prop:explicit-inc}
  यदि $\Gamma \Proves !A$ और $\lnot !A \in \Gamma$, तो $\Gamma$ असंगत है।
\end{prop}

\begin{proof}
  मान लें $\Gamma \Proves !A$ और $\lnot !A \in \Gamma$। तब
  $!B_1$, \dots, $!B_n \in \Gamma$ ऐसे हैं कि \
  \[
  \{\sFmla{\False}{!A}, \sFmla{\True}{!B_1}, \dots, \sFmla{\True}{!B_n}\}
  \]
  के लिए कोई बंद सत्य-वृक्ष है। मान्यता \sFmla{\False}{!A} को
  \sFmla{\True}{\lnot !A} से बदलें और मान्यताओं के बाद
  \sFmla{\False}{!A} पर \TRule{\True}{\lnot} लगाने का निष्कर्ष जोड़ें।
  प्रत्येक !!{sentence}, जो !!{tableau} में पुराने !!{tableau} की
  पंक्ति~$1$ के आधार पर उचित ठहराया गया था, अब पंक्ति~$n+1$ के आधार पर उचित
  ठहरता है। अतः यदि पुराना !!{tableau} बंद था, तो नया भी बंद है। चूँकि सभी मान्यताएँ~$\Gamma$ में
  हैं, यह दिखाता है कि $\Gamma$ असंगत है।
\end{proof}

\begin{prop}\ollabel{prop:provability-exhaustive}
  यदि $\Gamma \cup \{!A\}$ और $\Gamma \cup \{\lnot !A\}$ दोनों असंगत हैं,
  तो $\Gamma$ असंगत है।
\end{prop}

\begin{proof}
  यदि $!B_1$, \dots, $!B_n \in \Gamma$ और $!C_1$, \dots,
  $!C_m \in \Gamma$ ऐसे हैं कि
  \begin{align*}
    \{\sFmla{\True}{!A}, &
    \sFmla{\True}{!B_1}, \dots, \sFmla{\True}{!B_n}\} \text{ और}\\
    \{\sFmla{\True}{\lnot !A}, &
    \sFmla{\True}{!C_1}, \dots, \sFmla{\True}{!C_m}\}
  \end{align*}
  दोनों के लिए बंद !!{tableau}s हों, तो हम एक संयुक्त बंद !!{tableau} बना
  सकते हैं जो दिखाता है कि $\Gamma$ असंगत है। इसके लिए मान्यताओं के रूप में
  $\sFmla{\True}{!B_1}$, \dots, $\sFmla{\True}{!B_n}$ के साथ
  $\sFmla{\True}{!C_1}$, \dots, $\sFmla{\True}{!C_m}$ लेकर और फिर \Cut{}
  नियम लगाएँ। इससे दो शाखाएँ मिलती हैं: एक
  $\sFmla{\True}{!A}$ से और दूसरी $\sFmla{\False}{!A}$ से आरंभ होती है।
  
  बाईं ओर पहले !!{tableau} की मान्यताओं के नीचे का भाग जोड़ें। यहाँ नियम का
  प्रत्येक प्रयोग अब भी सही है, क्योंकि पहले !!{tableau} की प्रत्येक
  मान्यता, $\sFmla{\True}{!A}$ सहित, उपलब्ध है। अतः
  $\sFmla{\True}{!A}$ के नीचे की हर शाखा बंद हो जाती है।
  
  दाईं ओर दूसरे !!{tableau} की मान्यताओं के नीचे का भाग जोड़ें, पर
  $\TRule{\True}{\lnot}$ को $\sFmla{\True}{\lnot !A}$ पर लगाने से प्राप्त
  परिणाम हटा दें। $\TRule{\True}{\lnot}$ का
  $\sFmla{\True}{\lnot !A}$ पर प्रयोग करने का निष्कर्ष
  $\sFmla{\False}{!A}$ है, जो फिर भी
  उपलब्ध है, क्योंकि वह संयुक्त !!{tableau} की दाईं ओर \Cut{} नियम का
  निष्कर्ष है।

  यदि दूसरे सत्य-वृक्ष की कोई शाखा इसलिए बंद थी कि उसमें मान्यता
  $\sFmla{\True}{\lnot !A}$ (जो अब संयुक्त !!{tableau} में मान्यता के रूप
  में नहीं आती) के साथ $\sFmla{\False}{\lnot !A}$ भी था, तो
  $\TRule{\False}{\lnot}$ को $\sFmla{\False}{\lnot !A}$ पर लगाकर
  $\sFmla{\True}{!A}$ प्राप्त करें। अब संयुक्त !!{tableau} में संगत शाखा भी
  बंद हो जाती है, क्योंकि उसमें \Cut{} नियम का दायाँ निष्कर्ष
  $\sFmla{\False}{!A}$ है। यदि दूसरे !!{tableau} की कोई शाखा किसी अन्य
  कारण से बंद थी, तो संयुक्त !!{tableau} में संगत शाखा भी बंद है, क्योंकि
  $\sFmla{\True}{\lnot !A}$ के अतिरिक्त सभी !!{signed formula}s, जो पुराने
  दूसरे !!{tableau} की शाखा पर आते हैं, संयुक्त !!{tableau} की संगत शाखा पर
  भी आते हैं।
  \end{proof}

\end{document}
